\chapter{Códigos Computacionais em MATLAB}\label{anexo:codigos_matlab}

A reprodutibilidade integral dos cálculos numéricos de regime permanente e das rotinas de parametrização do ambiente Simscape SPS fundamenta-se na estruturação modular e vetorizada dos códigos-fonte desenvolvidos em MATLAB. A documentação das rotinas evidencia o emprego de técnicas computacionais eficientes que previnem descontinuidades numéricas e automatizam a geração dos resultados analíticos e gráficos requeridos.

% =========================================================================
% [DIRETRIZ DE REDAÇÃO]
% - Objetivo: Apresentar a documentação completa dos scripts computacionais desenvolvidos
%   em MATLAB para o cumprimento de sim_proj.pdf ("linhas de código").
% - Formulação: Disponibilizar as rotinas analíticas vetorizadas, o gerador da família V/f
%   e o script de setup de parâmetros e execução automatizada do Simscape SPS.
% - Figuras/Tabelas: Formatar as listagens de código através do ambiente lstlisting (style=matlabstyle).
% - Confronto: Garantir a reprodutibilidade dos resultados analíticos e das figuras do relatório.
% =========================================================================

\section{Rotina de Resolução Analítica Vetorizada e Varreduras Nominais}\label{sec:codigo_analitico}

A \Cref{lst:run_item_ABC} apresenta o código-fonte em MATLAB responsável pelo cálculo analítico do ponto nominal e pela geração vetorizada das curvas características de regime permanente.

\begin{lstlisting}[style=matlabstyle, caption={Rotina principal de calculo analitico e varredura de regime permanente.}, label={lst:run_item_ABC}]
%% ELE035 - TRABALHO FINAL: PARTE 1 (ITENS A, B, C, D, E)
% Resolucao Analitica Vetorizada do Circuito Equivalente IEEE
clear; clc; close all;

% Parametros nominais da maquina (WEG 3.7 kW, carcaca 100L)
P_nom = 3700;           % Potencia nominal [W]
V_LL = 220;             % Tensao de linha [V RMS]
V_ph = V_LL;            % Conexao Triangulo (Delta)
f_nom = 60;             % Frequencia [Hz]
p = 4;                  % Polos
n_nom = 1715;           % Rotacao nominal [rpm]
n_s = 120 * f_nom / p;  % Rotacao sincrona [rpm]
ws = 2 * pi * f_nom / (p/2); % Velocidade sincrona [rad/s]

% Parametros concentrados por fase em Triangulo
R1 = 1.189;             % Resistencia estatorica [ohm]
R2_prm = 1.180;         % Resistencia rotorica refletida [ohm]
X1 = 2.158;             % Reatancia dispersao estator [ohm]
X2_prm = 3.174;         % Reatancia dispersao rotor [ohm]
Xm = 49.601;            % Reatancia magnetizante [ohm]
B = 0.0426;             % Atrito viscoso da bancada [N*m*s/rad]
J = 0.701;              % Inercia total da bancada [kg*m^2]

% Item 1.B: Conjugado nominal util de placa
w_nom = n_nom * 2 * pi / 60;
Tn = P_nom / w_nom;     % Tn = 20.6020 N*m

% Item 1.D: Varredura puramente vetorizada de 0 a n_s
s = linspace(1e-4, 1.0, 2000);
n_vec = (1 - s) * n_s;
wm_vec = n_vec * 2 * pi / 60;

Z2_prm = R2_prm ./ s + 1j * X2_prm;
Zp = (1j * Xm .* Z2_prm) ./ (Z2_prm + 1j * Xm);
Zin = R1 + 1j * X1 + Zp;

I1_ph = V_ph ./ Zin;
I_line = sqrt(3) * abs(I1_ph);
I2_prm = I1_ph .* (1j * Xm ./ (Z2_prm + 1j * Xm));

Pag = 3 * (abs(I2_prm).^2) .* (R2_prm ./ s);
Pem = Pag .* (1 - s);
Tem = Pag / ws;

% Potencia ativa e fator de potencia
P_in = 3 * real(V_ph .* conj(I1_ph));
fp = cos(angle(Zin));
eta_em = Pem ./ P_in;
\end{lstlisting}

\section{Rotina de Simulação Dinâmica e Parametrização Simscape SPS}\label{sec:codigo_simulink_setup}

A \Cref{lst:setup_simulink} documenta o script de configuração de parâmetros do bloco \textit{Asynchronous Machine} com transformação de impedâncias Estrela-Triângulo.

\begin{lstlisting}[style=matlabstyle, caption={Script de parametrizacao e mapeamento Star-Delta para o Simscape SPS.}, label={lst:setup_simulink}]
%% ELE035 - TRABALHO FINAL: PARTE 2 (ITEM 2.A)
% Parametrizacao do Bloco Asynchronous Machine SI Units (Simscape SPS)
clear; clc;

% Dados nominais
Vn = 220;               % Tensao nominal de linha [V RMS]
Pn = 3700;              % Potencia nominal [VA equivalente]
fn = 60;                % Frequencia nominal [Hz]
p_pairs = 2;            % Pares de polos

% Transformacao de Impedancias Triangulo -> Estrela interna do bloco
R1_Y = 1.189 / 3;       % 0.39633 ohm
R2_Y = 1.180 / 3;       % 0.39333 ohm
Ll1_Y = 5.725e-3 / 3;   % 1.90833 mH
Ll2_Y = 8.420e-3 / 3;   % 2.80667 mH
Lm_Y = 131.571e-3 / 3;  % 43.8570 mH

% Parametros Mecanicos Invariantes
J_eq = 0.701;           % Momento de inercia acoplado [kg*m^2]
B_eq = 0.0426;          % Atrito viscoso [N*m*s/rad]
T_load_nominal = 20.602;% Degrau de carga nominal [N*m]
\end{lstlisting}
